% CONTAINDER --- Supplementary material, submitted as a SEPARATE file from the manuscript.
% Holds detail moved out of the body under the IJCIP length recommendation. Nothing here is a
% result the body does not already state.
%
% Sections are numbered SM.n rather than Sn: the body uses S0-S3 for the four mechanism-ablation
% arms, and "S1" would otherwise name both an ablation arm and a supplement section.
%
% NOTE: every numeral in this file is covered by check_numbers.py, which includes
% supplementary.tex in its targets. Moving a number out of the body must not move it out of the
% freeze gate.
\documentclass[review]{elsarticle}

\usepackage{amsmath,amssymb}
\usepackage{booktabs}
\PassOptionsToPackage{hyphens}{url}
\usepackage{url}

\journal{International Journal of Critical Infrastructure Protection}

\begin{document}

\begin{frontmatter}
\title{Supplementary material for\\
CONTAINDER: Bounding Credential-Compromise Impact in IEEE~2030.5 DER Systems}
\author{Huan Bui}
\ead{hbui11@charlotte.edu}
\affiliation{organization={University of North Carolina at Charlotte},
             city={Charlotte}, state={NC}, country={USA}}
\end{frontmatter}

\section*{SM.1 Standards quotations supporting the claim matrix}

The manuscript's standards table cites a free source for each claim. This section reproduces the
load-bearing text verbatim, so a reader can check the reading without obtaining any document that
must be purchased. All quotations are from Sandia National Laboratories, \emph{Recommendations for
Trust and Encryption in DER Interoperability Standards}, SAND2019-1490, whose full text is
publicly available from OSTI.

\paragraph{Certificate lifetime and revocation.} From \S4.1:

\begin{quote}
``The IEEE~2030.5/CSIP PKI is different than other PKI's in the use of non-expiring certificates
and the explicit prohibition of CRL and OCSP. Once issued, a device certificate has an indefinite
lifetime, so it is always valid. \dots\ If the device's private key is compromised, the device
certificate cannot be revoked. \dots\ In fact, IEEE~2030.5 prohibits their use. This means that
Certificate Authorities (CA) shall not maintain CRLs or run OSCP servers, and clients and servers
shall not check for CRLs or OCSP servers to verify certificate validity.''
\end{quote}

The prohibition is scoped, which is why the manuscript qualifies it. From \S6.4, quoting the
standard's own certificate-management statements:

\begin{quote}
``The phrase `Optional OCSP' means that the server device (and optionally, the client device) may
utilize Online Certificate Status Protocol as an additional mechanism to determine if a
certificate has been revoked. \textbf{OCSP may only be used to verify non-IEEE~2030.5
certificates.}''
\end{quote}

\paragraph{Local allow/deny listing, and its limits.} From \S4.1 and \S4.1.2:

\begin{quote}
``This does not preclude servers or clients from maintaining or obtaining their own list of
blacklisted or whitelisted devices if the operator prefers.''

``A blacklist is essentially a local CRL but without the process, rigor and verification behind
its issuance.''

``With fragmented blacklists and policies, it is possible for a stolen device certificate that was
disallowed in one region to be used in a different region to gain access because it was only
locally blacklisted.''
\end{quote}

\paragraph{A short-lived operational credential is permitted.} From \S6.4:

\begin{quote}
``Inverter manufacturers strictly following these guidelines are \textbf{not required to limit the
life of a device certificate} or to validate IEEE~2030.5 certificates using OCSP or CRLs.''

``(1 year certificate life is the usual recommended renewal interval).''
\end{quote}

This is why the manuscript claims no novelty for the short lifetime itself: it is recommended
practice in the analysis of record, and permitted rather than prohibited by the profile.

\paragraph{The adversary the design targets.} From \S6.4:

\begin{quote}
``If a manufacturer decides not to implement certificate revocation and replacement, for a
critical period an adversary could potentially extract a private key from a device and masquerade
as the legitimate device without other nodes being aware of the device compromise.''
\end{quote}

\paragraph{What no accessible source states.} Searched for and not found, in SAND2019-1490 or in
the CSIP Implementation Guide: any statement that adding an identity to a deny list terminates
that identity's \emph{established} mutual-TLS session; any statement that it cancels an
already-issued \texttt{DERControl}; and any mechanism by which an authorization can bound the
\emph{value} a control may take rather than which function may be invoked. The first two are the
premise of the manuscript's containment argument, and the third is why its numeric authorization
is placed above \texttt{FunctionSetAssignments} rather than inside it.

\section*{SM.2 The IEEE~9500 non-convergence, in full}

The manuscript states that the PNNL~9500-node system did not converge in our build and therefore
contributes no reported number. This section names the failure mode, so the absence is an
explained negative result rather than a silent omission.

The failure is in the \emph{base case}, with no attack and no PV of ours placed: the model's own
shipped driver, unaltered, returns \texttt{Converged = False} with an empty error description and
every bus per-unit magnitude reading zero --- the solution collapses rather than missing tolerance.
It reproduces at $1000$ iterations, under both the normal and Newton algorithms, and with feeder
controls enabled and disabled, so it is neither a control-loop oscillation nor an
iteration-budget shortfall. Our build is DSS C-API $0.14.5$ / DSS-Python $0.15.7$ /
OpenDSSDirect.py $0.9.4$ on macOS arm64.

\textbf{We did not resolve it and we do not attribute it}, having had no second OpenDSS build to
discriminate a defect on our side from one in the model. The consequence is that the 9500 system
supports nothing in the manuscript. Cross-feeder generality rests instead on the IEEE~8500-node
and IEEE~123-bus feeders, which every physical result in the paper uses, and which differ in
voltage level, topology, regulator structure and stiffness.

\section*{SM.3 Inventory of the runs behind the manuscript}

Table~\ref{tab:inventory} inventories every run, so no experiment's scale has to be inferred from
prose. Seeds are paired across arms within an experiment. ``Arms'' counts independent feeder
compiles; ``solves'' counts power-flow solutions, since an arm that searches an authorization set
or steps a horizon contains many.

\begin{table}[h]
\centering
\caption{Runs behind the manuscript. The calibration and the two post-hoc sweeps are labelled as
such; neither supports a hypothesis test.}
\label{tab:inventory}
\small
\begin{tabular}{@{}lrrrl@{}}
\toprule
Experiment & Seeds & Arms & Solves & Status \\
\midrule
Hosting-capacity calibration      & 5   & 560  & 1120  & calibration \\
Authorization-shape sweep         & 20  & 2520 & 18000 & confirmatory \\
Lifecycle horizon experiment      & 20  & 40   & 48000 & confirmatory \\
Attacker action-space scan        & 20  & 60   & 2760  & confirmatory \\
Reliance resolution sweep         & 20  & 160  & 1280  & post-hoc \\
Lifecycle parameter sensitivity   & --- & 60   & ---   & post-hoc, model only \\
Credential-service invariants     & 1   & 10   & ---   & confirmatory \\
Overhead microbenchmarks          & 500 iterations & 4 & --- & confirmatory \\
Availability fault injection      & 3 policies & 3 & --- & confirmatory \\
\bottomrule
\end{tabular}
\end{table}

The authorization-shape sweep is $2$ feeders $\times$ $3$ penetration rungs $\times$ $11$
authorization sets $\times$ $2$ control primitives $\times$ $20$ paired seeds, less the read-only
set, which has no fixed-setpoint realisation. The lifecycle experiment is $2$ operating states
$\times$ $20$ paired seeds $\times$ $19$ lifecycle arms plus a paired legitimate arm, each stepped
over a $60$-minute horizon.

\section*{SM.4 Result files and the claims they support}

Every numeral in the manuscript is checked against these files by \texttt{check\_numbers.py},
which fails if a reported number cannot be located in them.

\begin{table}[h]
\centering
\caption{Released result files and the manuscript claims each supports.}
\label{tab:files}
\small
\begin{tabular}{@{}ll@{}}
\toprule
File & Supports \\
\midrule
\texttt{hosting\_capacity.json}      & operating ladder, compliance criterion \\
\texttt{authz\_shape.json}           & authorization-set families, H1, H2, H3 \\
\texttt{shape\_contrasts.json}       & paired H1/H2/H3 estimates and intervals \\
\texttt{reliance\_resolution.json}   & intermediate rungs of the reliance curve \\
\texttt{lifecycle\_physical.json}    & horizon time series, all lifecycle arms \\
\texttt{lifecycle\_contrasts.json}   & paired lifecycle contrasts, Holm-corrected \\
\texttt{lifecycle\_sensitivity.json} & parameter sensitivity of the exposure window \\
\texttt{attackers.json}              & five attacker models, full action-space scan \\
\texttt{attacker\_detectors.json}    & dual-detector analysis, dual-evading optimum \\
\texttt{retention.json}              & compromise-class retention expectations \\
\texttt{credsvc.json}                & the ten security invariants \\
\texttt{hwkey.json}                  & the failed hardware-keystore attempt \\
\texttt{controlmode.json}            & feeder controls active under \texttt{static} \\
\texttt{availability.json}           & safe-mode availability under verifier outage \\
\bottomrule
\end{tabular}
\end{table}

\end{document}
